% !TEX encoding = UTF-8 Unicode
% !TEX TS-program = xelatex
% !TEX root = main.tex

\documentclass[xetex]{beamer}

%%%%%%%%%%%%%%%%%%%%%%%%%%%%%%%%%%%%%%%%%%%%%%%%%%%%%%%%%%%%%%%%%%%%
% Beamer settings

\mode<presentation>
{
  %\usetheme{Frankfurt}
  \useoutertheme[subsection=false]{smoothbars}
  \setbeamercovered{transparent}
  \setbeamertemplate{navigation symbols}{} %remove navigation symbols
  \setbeamercovered{invisible} %No transparent layers
}


\usepackage[many]{tcolorbox}
%\usepackage{tcolorbox}
\tcbuselibrary{skins,xparse}


\newcounter{ct}
\newcounter{contribution}

% put color to \boxed math command
\newcommand*{\boxcolor}{orange}
\makeatletter
\renewcommand{\boxed}[1]{\textcolor{\boxcolor}{%
\tikz[baseline={([yshift=-1ex]current bounding box.center)}] \node [rectangle, minimum width=1ex,rounded corners,draw] {\normalcolor\m@th$\displaystyle#1$};}}
 \makeatother
 \newcommand{\tboxed}[1]{\textcolor{\boxcolor}{%
\tikz[baseline={([yshift=-1ex]current bounding box.center)}] \node [rectangle, minimum width=1ex,rounded corners,draw] {\normalcolor#1};}}
 \makeatother

\usepackage{etoolbox}

\newtoggle{presentation}
\toggletrue{presentation}
%\togglefalse{presentation}

\newtoggle{euler}
\toggletrue{euler}
%\togglefalse{presentation}
 

% Loads configuration file

% !TEX encoding = UTF-8 Unicode
% !TEX TS-program = xelatex

%%%%%%%%%%%%%%%%%%%%%%%%%%%%%%%%%%%%%%%%%%%%%%%%%%%%%%%%%%%%%%%%%%%%
% AMS-like packages
\usepackage{amsmath,amsfonts,amssymb} % AMS fonts and symbols package
\usepackage{mathrsfs}

%%%%%%%%%%%%%%%%%%%%%%%%%%%%%%%%%%%%%%%%%%%%%%%%%%%%%%%%%%%%%%%%%%%%
% Typeface Settings

\usefonttheme{professionalfonts} %Necessary to use with Euler: This is a beamer setting!
\usepackage[small,euler-digits,euler-hat-accent]{eulervm} % Load BEFORE mathspec AND fontspec packages

\usepackage{mathspec}
		%\setmathrm[Mapping=tex-text,Numbers=OldStyle,Scale=MatchLowercase]{Palatino Linotype}

	\usepackage{fontspec}
		\defaultfontfeatures{Mapping=tex-text,Ligatures={TeX,Common},Scale=MatchUppercase,Color=textcolor}

		\setromanfont[
            Path = fonts/,
            Extension      = .ttf,
            UprightFont    = *,
            BoldFont       = *-Bold,
            ItalicFont     = *-Oblique,
            BoldItalicFont = *-BoldOblique]
			{Adobe_Garamond_Pro}
		
		\setsansfont[
					Extension 	   = .otf,
					UprightFont	   = *-regular,
					BoldFont	   = *-bold,
					ItalicFont	   = *-italic,
					BoldItalicFont = *-bolditalic]
					{texgyreheros}
		\setsansfont[
					Extension 	   = .otf,
					UprightFont	   = *-regular,
					BoldFont	   = *-bold,
					ItalicFont	   = *-italic,
					BoldItalicFont = *-bolditalic]
					{texgyreheros}
		% Do not add "Numbers=OldStyle" for the options of the fonts below
		% \setmonofont[UprightFont 	= Inconsolata-Regular.ttf,
		% 			 BoldFont   	= Inconsolata-Bold.ttf]
		% 			{Inconsolata-Regular.ttf}
		
		% \newfontfamily\thinfont[Numbers=OldStyle,ItalicFont=FuturaLT-LightOblique.ttf]{FuturaLT-Light.ttf}

		% \newfontfamily{\normal}[Mapping=tex-text,Scale=MatchLowercase]{Futura LT Medium}
		% \newfontfamily{\light}[Mapping=tex-text,Scale=MatchLowercase]{Futura LT Light}
		
% Definition of the \newhtought, \allcaps, \smallcaps and \textsc

\usepackage{textcase}

\newskip\tufteskipamount
\tufteskipamount=1.0\baselineskip plus 0.5ex minus 0.2ex

\newcommand{\tuftebreak}{\par\ifdim\lastskip<\tufteskipamount
  \removelastskip\penalty-100\tufteskip\fi}

\newcommand{\tufteskip}{\vspace\tufteskipamount}

\def\smallcapsspacing{\@tufte@warning{Proper spacing of small-caps letters has not been set up.}}

\DeclareTextFontCommand{\textsmallcaps}{\scshape}
\def\allcapsspacing{\@tufte@warning{Proper spacing of ALL-CAPS letters has not been set up.}}
\def\smallcapsspacing{\@tufte@warning{Proper spacing of small-caps letters has not been set up.}}
\newcommand{\allcaps}[1]{\allcapsspacing{\MakeTextUppercase{#1}}}
\newcommand{\smallcaps}[1]{\smallcapsspacing{\MakeTextLowercase{#1}}}

\newcommand{\textls}[2][5]{%
    \begingroup\addfontfeatures{LetterSpace=#1}#2\endgroup
  }
  
  \providecommand\newthought[1]{%
   \tuftebreak
   \noindent\textsc{#1}%
}
\renewcommand{\allcapsspacing}[1]{\textls[15]{#1}}
 \renewcommand{\smallcapsspacing}[1]{\textls[10]{#1}}
 \renewcommand{\allcaps}[1]{\textls[15]{\MakeTextUppercase{#1}}}
 \renewcommand{\smallcaps}[1]{\smallcapsspacing{\scshape\MakeTextLowercase{#1}}}
 \renewcommand{\textsc}[1]{\smallcapsspacing{\textsmallcaps{#1}}}

% End of the definition of the \newthought command ---------

%%%%%%%%%%%%%%%%%%%%%%%%%%%%%%%%%%%%%%%%%%%%%%%%%%%%%%%%%%%%%%%%%%%%
% Language Settings

% Needs to be loaded locally and after the fontspec package

\usepackage{polyglossia}
\setdefaultlanguage[variant=british]{english}
\setotherlanguages{french,portuges,italian}

\usepackage{csquotes}


%%%%%%%%%%%%%%%%%%%%%%%%%%%%%%%%%%%%%%%%%%%%%%%%%%%%%%%%%%%%%%%%%%%%
% Graphic settings


\usepackage{graphicx,xcolor,epstopdf}
\usepackage{color}
\DeclareGraphicsRule{.tif}{png}{.png}{`convert #1 `dirname #1`/`basename #1 .tif`.png}

% Color definitions
\definecolor{canary}{HTML}{FCF0AD}

\definecolor{gindigo}{HTML}{2196F3}


%Blue

\definecolor{gblue500}{HTML}{2196F3}
\definecolor{gblue50}{HTML}{E3F2FD}
\definecolor{gblue100}{HTML}{BBDEFB}
\definecolor{gblue200}{HTML}{90CAF9}
\definecolor{gblue300}{HTML}{64B5F6}
\definecolor{gblue400}{HTML}{42A5F5}
\definecolor{gblue600}{HTML}{1E88E5}
\definecolor{gblue700}{HTML}{1976D2}
\definecolor{gblue800}{HTML}{1565C0}
\definecolor{gbluetext}{HTML}{1565C0}
\definecolor{gblue900}{HTML}{0D47A1}

%Green


\definecolor{ggreen500}{HTML}{4CAF50}
\definecolor{ggreen50}{HTML}{E8F5E9}
\definecolor{ggreen100}{HTML}{C8E6C9}
\definecolor{ggreen200}{HTML}{A5D6A7}
\definecolor{ggreen300}{HTML}{81C784}
\definecolor{ggreen400}{HTML}{66BB6A}
\definecolor{ggreen500}{HTML}{4CAF50}
\definecolor{ggreen600}{HTML}{43A047}
\definecolor{ggreen700}{HTML}{388E3C}
\definecolor{ggreen800}{HTML}{2E7D32}
\definecolor{ggreen900}{HTML}{1B5E20}

%Red


\definecolor{ggrey500}{HTML}{9E9E9E}
\definecolor{ggrey50}{HTML}{FAFAFA}
\definecolor{ggrey100}{HTML}{F5F5F5}
\definecolor{ggrey200}{HTML}{EEEEEE}
\definecolor{ggrey300}{HTML}{E0E0E0}
\definecolor{ggrey400}{HTML}{BDBDBD}
\definecolor{ggrey500}{HTML}{9E9E9E}
\definecolor{ggrey600}{HTML}{757575}
\definecolor{ggrey700}{HTML}{616161}
\definecolor{ggrey800}{HTML}{424242}
\definecolor{ggrey900}{HTML}{212121}

%Grey


\definecolor{gred500}{HTML}{F44336}
\definecolor{gred50}{HTML}{FFEBEE}
\definecolor{gred100}{HTML}{FFCDD2}
\definecolor{gred200}{HTML}{EF9A9A}
\definecolor{gred300}{HTML}{E57373}
\definecolor{gred400}{HTML}{EF5350}
\definecolor{gred500}{HTML}{F44336}
\definecolor{gred600}{HTML}{E53935}
\definecolor{gred700}{HTML}{D32F2F}
\definecolor{gred800}{HTML}{C62828}
\definecolor{gredtext}{HTML}{C62828}
\definecolor{gred900}{HTML}{B71C1C}

%Yellow

\definecolor{gyellow500}{HTML}{FFEB3B}
\definecolor{gyellow50}{HTML}{FFFDE7}
\definecolor{gyellow100}{HTML}{FFF9C4}
\definecolor{gyellow200}{HTML}{FFF59D}
\definecolor{gyellow300}{HTML}{FFF176}
\definecolor{gyellow400}{HTML}{FFEE58}
\definecolor{gyellow600}{HTML}{FDD835}
\definecolor{gyellow700}{HTML}{FBC02D}
\definecolor{gyellow800}{HTML}{F9A825}
\definecolor{gyellow900}{HTML}{F57F17}

% Orange

\definecolor{gorange500}{HTML}{FF9800}

% Postit color

\definecolor{postitcolor}{HTML}{FCF0AD}

\usepackage{pdfrender}
\usepackage[beamer]{hf-tikz}
\usepackage{tikz}
\usetikzlibrary{arrows,shapes,shadings}
\usetikzlibrary{calc}

\usepackage{multimedia,animate, colortbl}

\usepackage{tcolorbox} 
\tcbuselibrary{skins,theorems}
\usepackage{varwidth} % To change blocks width

%%%%%%%%%%%%%%%%%%
% OTHER PACKAGES %
%%%%%%%%%%%%%%%%%%

% Separates counting from actual presentation from appendix
\usepackage{appendixnumberbeamer}

% Reduce left margin
\usepackage{changepage}

% Conditional package
\usepackage{ifthen,forloop}

\usepackage{multicol}
\setlength{\columnseprule}{0.4pt}

% Bla bla bla generator
\usepackage{blindtext}

%%%%%%%%%%%%%%
% Water mark %
%%%%%%%%%%%%%%

% different for odd and even pages
%\usepackage{draftwatermark}
%
%\makeatletter
%\renewcommand\sc@wm@print[1]{% redefine positioning of mark (-1in to 330pt)
% \if@sc@wm@stamp
%          \setbox\@tempboxa\vbox to \z@{%
%          \ifodd\c@page
%            \vskip 2em  \moveright 8em
%           \else
%            \vskip 2em  \moveright 12.5em
%           \fi \vbox{%
%              \hbox to \z@{%
%                #1\hss}}\vss}
%          \dp\@tempboxa\z@
%          \box\@tempboxa
%          \fi}
%\makeatother
%
%\makeatletter
%\SetWatermarkText{%
%    \ifodd\c@page
%        \begin{tikzpicture}\node[opacity=.1]		{\includegraphics[angle=315,scale=0.0625]{./imgs/MIT}};				\end{tikzpicture}
%    \else
%    	\begin{tikzpicture}\node[opacity=.1]		{\includegraphics[angle=315,scale=1]{./imgs/Usyd_new_logo}};
%   		\end{tikzpicture}
%    \fi
%}
%\makeatother
%
%\setbeamertemplate{background canvas}{}

%%%%%%%%%%%%%%%%%%%%
% HYPERREF PACKAGE %
%%%%%%%%%%%%%%%%%%%%
\usepackage{hyperref}% backref linktocpage pagebackref

\hypersetup{
% Uncomment the line below to remove all links (to references, figures, tables, etc)
%draft, 
colorlinks=true, linktocpage=true, pdfstartpage=1, pdfstartview=FitV,
% Uncomment the line below if you want to have black links (e.g. for printing black and white)
colorlinks=false, linktocpage=false, pdfborder={0 0 0}, pdfstartpage=1, pdfstartview=FitV, 
%
% PDF file meta-information
%
pdftitle={},
pdfauthor={Humberto Stein Shiromoto},
pdfsubject={},
pdfkeywords={},
pdfcreator={pdfLaTeX},
pdfproducer={}
}

\setbeamercolor{bibliography entry author}{fg=ggreentext}
\setbeamercolor{bibliography entry title}{fg=black}
\setbeamercolor{bibliography entry location}{fg=gredtext}
\setbeamercolor{bibliography entry Publisher
}{fg=gredtext}
\setbeamercolor{bibliography entry note
}{fg=black}
\setbeamercolor{bibliography item}{fg=gblue700}

%%%%%%%%%%%%%%%%%%%%%%%%%%%%%%%%%%%%%%%
% BEAMER ITEMIZE AND ENUMERATE STYLES %
%%%%%%%%%%%%%%%%%%%%%%%%%%%%%%%%%%%%%%%

% \usepackage{enumerate} %No need to be loaded with Beamer

\setbeamertemplate{itemize item}[square]
\setbeamercolor{itemize item}{fg=gblue700}

 % if you want a circle
\setbeamertemplate{itemize subitem}[circle]
\setbeamercolor{itemize subitem}{fg=gblue700}

\setbeamertemplate{enumerate item}[square]
\setbeamercolor{enumerate item}{fg=gblue700,bg=gblue700}

\setbeamercolor{local structure}{fg=gblue700}

\usepackage{verbatim}


% Removes shadow from title page box
\setbeamertemplate{title page}[default][colsep=-4bp,rounded=true]

%%%%%%%%%%%%%%
% Frametitle %
%%%%%%%%%%%%%%

% Removes shadow from frametitle
\setbeamertemplate{frametitle}[default][colsep=-4bp,rounded=false,shadow=false]


\setbeamercolor{frametitle}{fg=white,bg=gblue700}


%%%%%%%%%%%%
% Headline %
%%%%%%%%%%%%
%\setbeamerfont{headline}{series=\bfseries}
\setbeamercolor{section in head/foot}{fg=white, bg=black}
\setbeamercolor{upper separation line head}{bg=white}
\setbeamertemplate{section in head/foot shaded}{\color{white!50!black}\insertsectionhead}

\makeatletter
\setbeamertemplate{headline}
{%
  \pgfuseshading{beamer@barshade}%
  \ifbeamer@sb@subsection%
    \vskip-9.75ex%
  \else%
    \vskip-7ex%
  \fi%
  \begin{beamercolorbox}[ignorebg,ht=2.25ex,dp=3.75ex]{section in head/foot}
    \insertnavigation{\paperwidth}
  \end{beamercolorbox}%
  \ifbeamer@sb@subsection%
    \begin{beamercolorbox}[ignorebg,ht=2.125ex,dp=1.125ex,%
      leftskip=.3cm,rightskip=.3cm plus1fil]{subsection in head/foot}
      \usebeamerfont{subsection in head/foot}\insertsubsectionhead
    \end{beamercolorbox}%
  \fi%
 \begin{beamercolorbox}[colsep=1.5pt,ht=0.75ex]{upper separation line head}
 \end{beamercolorbox}
}%
\makeatother


%%%%%%%%%%%%
% Footline %
%%%%%%%%%%%%
%\setbeamertemplate{footline}[frame number]
\setbeamercolor{footline}{fg=white,bg=black}

\makeatletter
\defbeamertemplate*{footline}{my footline}
{
   \ifnum \insertpagenumber=1
      \leavevmode%
      \hbox{%
      \begin{beamercolorbox}[wd=\paperwidth,ht=2.25ex,dp=1ex,center]{}%
        % empty environment to raise height
      \end{beamercolorbox}}%
      \vskip0pt%
    \else
  \leavevmode%
  \hbox{%
  \begin{beamercolorbox}[wd=.5\paperwidth,ht=2.25ex,dp=1ex,center]{footline}%
    \usebeamerfont{author in head/foot}\insertshortauthor
  \end{beamercolorbox}%
  \begin{beamercolorbox}[wd=.5\paperwidth,ht=2.25ex,dp=1ex,center]{footline}%
    \usebeamerfont{title in head/foot}\insertshorttitle\hfill(\insertframenumber{} / \inserttotalframenumber)
  \end{beamercolorbox}
  }%
  \vskip0pt%
  \fi
}
\makeatother


\setbeamertemplate{footline}[my footline]

%%%%%%%%%%%%%%%%%%%%%
% SECTION-TITLEPAGE %
%%%%%%%%%%%%%%%%%%%%%

\AtBeginSection[]{%
    {
    \setbeamercolor{upper separation line head}{bg=gblue700}
    \begin{frame}
    \begin{tikzpicture}[remember picture,overlay]
    \fill[gblue700]
        ([yshift=0pt]current page.west) rectangle ([yshift=-\headheight] current page.north east);
    \node[anchor=center]
        at ([yshift=10pt,xshift=0pt]current page.center) 
        {\parbox[t]{\textwidth}{\centering%
            \usebeamerfont{author}\textcolor{white}{%
                    \textpdfrender{
                        TextRenderingMode=FillStroke,
                        FillColor=white,
                        LineWidth=.1ex,
                    }{\insertsection}}}};
    \end{tikzpicture}
    \end{frame}}
    \addtocounter{framenumber}{-1}
}

%%%%%%%%%%%%%%%%%%%%%
% TABLE OF CONTENTS %
%%%%%%%%%%%%%%%%%%%%%

%Changing numbering in TOC
%\setbeamertemplate{sections/subsections in toc}[square]
%\setbeamertemplate{section in toc}{\inserttocsection}
%\setbeamerfont{section in toc}{series=\bfseries}
%
%% TOC must be loaded before \begin{document}
%\AtBeginSection[]{
%	\begin{frame}<beamer>
%		%\frametitle{Overview}
%		\tableofcontents[currentsection,hideallsubsections,subsubsectionstyle=hide]
%	\end{frame}
%	\addtocounter{framenumber}{-1} 
%}


% !TEX encoding = UTF-8 Unicode

% Mahbb commands

\newcommand{\bb}[1]{\mathbb{#1}}

\newcommand{\bbA}{\mathbb{A}}
\newcommand{\bbB}{\mathbb{B}}
\newcommand{\bbC}{\mathbb{C}}
\newcommand{\bbD}{\mathbb{D}}
\newcommand{\bbE}{\mathbb{E}}
\newcommand{\bbF}{\mathbb{F}}
\newcommand{\bbG}{\mathbb{G}}
\newcommand{\bbI}{\mathbb{I}}
\newcommand{\bbJ}{\mathbb{J}}
\newcommand{\bbK}{\mathbb{K}}
\newcommand{\bbL}{\mathbb{L}}
\newcommand{\bbM}{\mathbb{M}}
\newcommand{\bbN}{\mathbb{N}}
\newcommand{\bbO}{\mathbb{O}}
\newcommand{\bbP}{\mathbb{P}}
\newcommand{\bbQ}{\mathbb{Q}}
\newcommand{\R}{\mathbb{R}}
\newcommand{\Rs}{\mathbb{R}_{>0}}
\newcommand{\Ras}{\mathbb{R}_{\geq0}}
\newcommand{\bbS}{\mathbb{S}}
\newcommand{\bbT}{\mathbb{T}}
\newcommand{\bbU}{\mathbb{U}}
\newcommand{\bbV}{\mathbb{V}}
\newcommand{\bbX}{\mathbb{X}}
\newcommand{\bbY}{\mathbb{Y}}
\newcommand{\bbW}{\mathbb{W}}
\newcommand{\bbZ}{\mathbb{Z}}
\newcommand{\Zas}{\mathbb{Z}_{\geq0}}
\newcommand{\Zs}{\mathbb{Z}_{>0}}

% Mathbold commands

\iftoggle{euler}{%

\newcommand{\bA}{\mathbold{A}}
\newcommand{\bB}{\mathbold{B}}
\newcommand{\bC}{\mathbold{C}}
\newcommand{\bD}{\mathbold{D}}
\newcommand{\bE}{\mathbold{E}}
\newcommand{\bF}{\mathbold{F}}
\newcommand{\bG}{\mathbold{G}}
\newcommand{\bH}{\mathbold{H}}
\newcommand{\bI}{\mathbold{I}}
\newcommand{\bJ}{\mathbold{J}}
\newcommand{\bK}{\mathbold{K}}
\newcommand{\bL}{\mathbold{L}}
\newcommand{\bM}{\mathbold{M}}
\newcommand{\bN}{\mathbold{N}}
\newcommand{\bO}{\mathbold{O}}
\newcommand{\bP}{\mathbold{P}}
\newcommand{\bQ}{\mathbold{Q}}
\newcommand{\bR}{\mathbold{R}}
\newcommand{\bS}{\mathbold{S}}
\newcommand{\bT}{\mathbold{T}}
\newcommand{\bU}{\mathbold{U}}
\newcommand{\bV}{\mathbold{V}}
\newcommand{\bX}{\mathbold{X}}
\newcommand{\bY}{\mathbold{Y}}
\newcommand{\bW}{\mathbold{W}}
\newcommand{\bZ}{\mathbold{Z}}

}{
\newcommand{\bd}[1]{\mathbf{#1}}

\newcommand{\bA}{\mathbf{A}}
\newcommand{\bB}{\mathbf{B}}
\newcommand{\bC}{\mathbf{C}}
\newcommand{\bD}{\mathbf{D}}
\newcommand{\bE}{\mathbf{E}}
\newcommand{\bF}{\mathbf{F}}
\newcommand{\bG}{\mathbf{G}}
\newcommand{\bH}{\mathbf{H}}
\newcommand{\bI}{\mathbf{I}}
\newcommand{\bJ}{\mathbf{J}}
\newcommand{\bK}{\mathbf{K}}
\newcommand{\bL}{\mathbf{L}}
\newcommand{\bM}{\mathbf{M}}
\newcommand{\bN}{\mathbf{N}}
\newcommand{\bO}{\mathbf{O}}
\newcommand{\bP}{\mathbf{P}}
\newcommand{\bQ}{\mathbf{Q}}
\newcommand{\bR}{\mathbf{R}}
\newcommand{\bS}{\mathbf{S}}
\newcommand{\bT}{\mathbf{T}}
\newcommand{\bU}{\mathbf{U}}
\newcommand{\bV}{\mathbf{V}}
\newcommand{\bX}{\mathbf{X}}
\newcommand{\bY}{\mathbf{Y}}
\newcommand{\bW}{\mathbf{W}}
\newcommand{\bZ}{\mathbf{Z}}

}

% Mathcal commands

\newcommand{\cA}{\mathcal{A}}
\newcommand{\cB}{\mathcal{B}}
\newcommand{\cC}{\mathcal{C}}
\newcommand{\cD}{\mathcal{D}}
\newcommand{\cE}{\mathcal{E}}
\newcommand{\cF}{\mathcal{F}}
\newcommand{\cH}{\mathcal{H}}
\newcommand{\cI}{\mathcal{I}}
\newcommand{\cJ}{\mathcal{J}}
\newcommand{\cK}{\mathcal{K}}
\newcommand{\cL}{\mathcal{L}}
\newcommand{\cM}{\mathcal{M}}
\newcommand{\cN}{\mathcal{N}}
\newcommand{\cO}{\mathcal{O}}
\newcommand{\cP}{\mathcal{P}}
\newcommand{\cQ}{\mathcal{Q}}
\newcommand{\cR}{\mathcal{R}}
\newcommand{\cS}{\mathcal{S}}
\newcommand{\cT}{\mathcal{T}}
\newcommand{\cU}{\mathcal{U}}
\newcommand{\cV}{\mathcal{V}}
\newcommand{\cX}{\mathcal{X}}
\newcommand{\cY}{\mathcal{Y}}
\newcommand{\cW}{\mathcal{W}}
\newcommand{\cZ}{\mathcal{Z}}

% Mathscr commands

\newcommand{\scr}[1]{\mathscr{#1}}

\newcommand{\scrA}{\mathscr{A}}
\newcommand{\scrB}{\mathscr{B}}
\newcommand{\scrC}{\mathscr{C}}
\newcommand{\scrD}{\mathscr{D}}
\newcommand{\scrE}{\mathscr{E}}
\newcommand{\scrG}{\mathscr{G}}
\newcommand{\scrI}{\mathscr{I}}
\newcommand{\scrJ}{\mathscr{J}}
\newcommand{\scrK}{\mathscr{K}}
\newcommand{\scrL}{\mathscr{L}}
\newcommand{\scrM}{\mathscr{M}}
\newcommand{\scrN}{\mathscr{N}}
\newcommand{\scrO}{\mathscr{O}}
\newcommand{\scrP}{\mathscr{P}}
\newcommand{\scrQ}{\mathscr{Q}}
\newcommand{\scrR}{\mathscr{R}}
\newcommand{\scrS}{\mathscr{S}}
\newcommand{\scrT}{\mathscr{T}}
\newcommand{\scrU}{\mathscr{U}}
\newcommand{\scrV}{\mathscr{V}}
\newcommand{\scrX}{\mathscr{X}}
\newcommand{\scrY}{\mathscr{Y}}
\newcommand{\scrW}{\mathscr{W}}
\newcommand{\scrZ}{\mathscr{Z}}

% Mathds commands

\newcommand{\ds}[1]{\mathds{#1}}

\newcommand{\dsA}{\mathds{A}}
\newcommand{\dsB}{\mathds{B}}
\newcommand{\dsC}{\mathds{C}}
\newcommand{\dsD}{\mathds{D}}
\newcommand{\dsE}{\mathds{E}}
\newcommand{\dsI}{\mathds{I}}
\newcommand{\dsJ}{\mathds{J}}
\newcommand{\dsK}{\mathds{K}}
\newcommand{\dsL}{\mathds{L}}
\newcommand{\dsM}{\mathds{M}}
\newcommand{\dsN}{\mathds{N}}
\newcommand{\dsO}{\mathds{O}}
\newcommand{\dsP}{\mathds{P}}
\newcommand{\dsQ}{\mathds{Q}}
\newcommand{\dsR}{\mathds{R}}
\newcommand{\dsS}{\mathds{S}}
\newcommand{\dsT}{\mathds{T}}
\newcommand{\dsU}{\mathds{U}}
\newcommand{\dsV}{\mathds{V}}
\newcommand{\dsX}{\mathds{X}}
\newcommand{\dsY}{\mathds{Y}}
\newcommand{\dsW}{\mathds{W}}
\newcommand{\dsZ}{\mathds{Z}}

% Mathfrak commands

\newcommand{\kB}{\mathfrak{B}}
\newcommand{\kF}{\mathfrak{F}}
\newcommand{\kM}{\mathfrak{M}}
\newcommand{\kS}{\mathfrak{S}}
\newcommand{\kT}{\mathfrak{T}}

% Ohter Math commands

\newcommand{\mrm}[1]{\mathrm{#1}}
\newcommand{\mbf}[1]{\mathbf{#1}}
\newcommand{\mf}[1]{\mathfrak{#1}}


% !TEX encoding = UTF-8 Unicode

% To redeclare a MathOperator: begin
% Not compatible with arXiv !
\makeatletter
\newcommand\RedeclareMathOperator{%
  \@ifstar{\def\rmo@s{m}\rmo@redeclare}{\def\rmo@s{o}\rmo@redeclare}%
}
% this is taken from \renew@command
\newcommand\rmo@redeclare[2]{%
  \begingroup \escapechar\m@ne\xdef\@gtempa{{\string#1}}\endgroup
  \expandafter\@ifundefined\@gtempa
     {\@latex@error{\noexpand#1undefined}\@ehc}%
     \relax
  \expandafter\rmo@declmathop\rmo@s{#1}{#2}}
% This is just \@declmathop without \@ifdefinable
\newcommand\rmo@declmathop[3]{%
  \DeclareRobustCommand{#2}{\qopname\newmcodes@#1{#3}}%
}
\@onlypreamble\RedeclareMathOperator
\makeatother
% end %

% Commands

\newcommand{\dx}{\dot{x}}
\newcommand{\dy}{\dot{y}}
\newcommand{\xp}{x^+}
\newcommand{\lloc}{\mathcal{L}_\mathrm{loc}}
\newcommand{\cl}{\mathtt{cl}}
\newcommand{\dom}{\mathtt{dom}}
\newcommand{\domp}{\mathtt{dom}_{\geq0}}
\newcommand{\im}{\mathtt{im}}
\newcommand{\OSC}{\mathcal{OSC}}
\newcommand{\AC}{\mathcal{AC}}
\newcommand{\supp}{\mathtt{supp}}
\newcommand{\MAF}{\mathtt{MAF}}
\newcommand{\id}{\mathtt{id}}
\newcommand{\conv}{\mathtt{conv}}
\newcommand{\sign}{\mathtt{sign}}
\newcommand{\Liploc}{\cL\mathpzc{ip}_\mathrm{loc}}
\newcommand{\Ploc}{\mathcal{P}_{\mathrm{loc}}}
\renewcommand{\ae}{\mathrm{for\ a.e.\ }}
\RedeclareMathOperator{\Re}{\mathtt{Re}}
\RedeclareMathOperator{\Im}{\mathtt{Im}}
\newcommand*{\qedf}{\hfill\ensuremath{\blacksquare}}
\newcommand*{\qede}{\hfill\ensuremath{\blacksquare}}

\newcommand{\pfrac}[2]{\raisebox{1pt}{\small \ensuremath{#1}}$/$\raisebox{-2pt}{\small \ensuremath{#2}}}

% Commands with tt


% Operators

\DeclareMathOperator{\range}{\mathtt{range}}
\DeclareMathOperator{\wnull}{\mathtt{null}}
\DeclareMathOperator{\diag}{\mathtt{diag}}
\DeclareMathOperator{\dist}{\mathtt{dist}}
\DeclareMathOperator{\diverg}{\mathtt{div}}
\RedeclareMathOperator{\div}{\mathtt{div}}
\DeclareMathOperator{\sat}{\mathtt{sat}}
\DeclareMathOperator{\real}{\mathtt{real}}
\DeclareMathOperator{\rank}{\mathtt{rank}}

\DeclareMathOperator*{\argmin}{\mathtt{argmin}} %Note: the * that follows \DeclareMathOperator sets the underscored option (the argument) underneath the operator like the \lim operator.

\DeclareMathOperator{\interior}{\mathtt{int}}
\DeclareMathOperator{\ext}{\mathtt{ext}}
\DeclareMathOperator{\dir}{\mathrm{dir}}
\DeclareMathOperator{\gph}{\mathrm{gph}}
\DeclareMathOperator{\card}{\mathtt{card}}
\DeclareMathOperator{\diam}{\mathtt{diam}}
\DeclareMathOperator{\co}{\mathtt{co}}
\DeclareMathOperator{\mes}{\mathtt{mes}}
\DeclareMathOperator{\vol}{\mathtt{vol}}
\DeclareMathOperator{\proj}{\mathtt{proj}}
\DeclareMathOperator{\grad}{\mathtt{grad}}
\DeclareMathOperator*{\esssup}{\mathtt{ess\ sup}}

% Colors

\newcommand{\black}{\color{black}}
\newcommand{\blue}{\color{blue}}
\newcommand{\red}{\color{red}}
\newcommand{\new}{\color{black}}
\renewcommand{\stop}{\color{black}}

%\newcommand{\dif}[1]{\dot{\overgroup{\delta #1}}}
%\ifxetex
%	\newcommand{\widedot}[1]{\dot{\overgroup{#1}}}
%\else
%	\renewcommand{\widedot}[1]{\dot{\overgroup{#1}}}
%\fi

\definecolor{white}{HTML}{FFFFFF}
\newcommand{\white}{\color{white}}
% !TEX encoding = UTF-8 Unicode

% \usepackage{ntheorem} % Do not use this package
\usepackage{amsthm}
\usepackage{thmtools}

\iftoggle{presentation}{%

%%%%%%%%%%%%%%%%%%%%%%%%%%%%%%%%%%%%%%%%%%%%%%%%%%%%%%%%%%%%%%%%%%%%%
%																	%
% Block environments												%
%																	%
%%%%%%%%%%%%%%%%%%%%%%%%%%%%%%%%%%%%%%%%%%%%%%%%%%%%%%%%%%%%%%%%%%%%%

% \setbeamertemplate{blocks}[default]
\setbeamertemplate{blocks}[rounded][shadow=true]
%\setbeamerfont{block title}{series=\bfseries}

% Normal block environment settings
\setbeamercolor{structure}{bg=gblue700,fg=white}
\setbeamercolor*{block title}{fg=gblue700,
bg=white}
\setbeamercolor*{block body}{fg=black,
bg=white}

\makeatletter
\pgfdeclareverticalshading{bmb@transition}{200cm}{%
  color(0pt)=(gblue700); color(2pt)=(gblue700)}
\makeatother


% Example block settings
\setbeamercolor*{block title example}{fg=white,
bg=ggreen700}
\setbeamercolor*{block body example}{fg=black,
bg=ggreen100}
\setbeamercolor{example text}{fg=black}

% Alert block settings
\setbeamercolor*{block title alerted}{fg=white,
bg=gred500}
\setbeamercolor*{block body alerted}{fg=black,
bg=gred100}
\setbeamercolor{alerted text}{fg=black}

\newcommand*\redblock{enhanced,colback=gred50,
boxrule=0.4pt,sharp corners,
colframe=gred900,fonttitle=\bfseries}

\newcommand*\postitblock{enhanced,colback=postitcolor,
boxrule=0.4pt,sharp corners,
colframe=postitcolor,fonttitle=\bfseries}

\newcommand*\whiteblock{enhanced,colback=white,
boxrule=0.4pt,sharp corners,
colframe=ggrey300,fonttitle=\bfseries}

\tcbset{enhanced,colback=gblue50,
boxrule=0.4pt,sharp corners,
colframe=gblue700,fonttitle=\bfseries}

%% Testing blocks

\newtcolorbox{dblock}[2][]{
  beamer,
  colback=white,
  colframe=white,
  coltitle=gblue700,
  fonttitle=\normalfont\large,
  overlay={
 %   \draw[gblue700, thick] 
    \draw[gblue700, ultra thick] 
    ([xshift=6pt,yshift=-0.5\pgflinewidth]title.south west) --
    ([xshift=-6pt,yshift=-0.5\pgflinewidth]title.south east);},
  left=1mm,
  right=1mm,
  top=1mm,
  bottom=1mm,
  middle=1mm,
  fuzzy shadow={0.6ex}{-0.6ex}{1pt}{0.5pt}{black!80!white},
  #1,title=#2}
\makeatletter
\tikzset{beamer@color/.style={top color=white,bottom color=white}}
\makeatother

\newtcolorbox{tblock}[2][]{
  beamer,
  colback=white,
  colframe=white,
  coltitle=ggreen700,
  fonttitle=\normalfont\large,
  overlay={
%    \draw[ggreen700, thick] 
    \draw[ggreen700, ultra thick] 
    ([xshift=6pt,yshift=-0.5\pgflinewidth]title.south west) --
    ([xshift=-6pt,yshift=-0.5\pgflinewidth]title.south east);},
  left=1mm,
  right=1mm,
  top=1mm,
  bottom=1mm,
  middle=1mm,
  fuzzy shadow={0.6ex}{-0.6ex}{1pt}{0.5pt}{black!80!white},
  #1,title=#2}
\makeatletter
\tikzset{beamer@color/.style={top color=white,bottom color=white}}
\makeatother

\newtcolorbox{oblock}[2][]{
  beamer,
  colback=white,
  colframe=white,
  coltitle=gorange700,
  fonttitle=\normalfont\large,
  overlay={
%    \draw[ggreen700, thick] 
    \draw[gorange700, ultra thick] 
    ([xshift=6pt,yshift=-0.5\pgflinewidth]title.south west) --
    ([xshift=-6pt,yshift=-0.5\pgflinewidth]title.south east);},
  left=1mm,
  right=1mm,
  top=1mm,
  bottom=1mm,
  middle=1mm,
  fuzzy shadow={0.6ex}{-0.6ex}{1pt}{0.5pt}{black!80!white},
  #1,title=#2}
\makeatletter
\tikzset{beamer@color/.style={top color=white,bottom color=white}}
\makeatother

% For custom item with the color of the theorem block: currently not working
\newenvironment{tlistenv}{\only{\setbeamercolor{local structure}{fg=gorange700}}}{}

\newtcolorbox{rblock}[2][]{
  beamer,
  colback=white,
  colframe=white,
  coltitle=black,
  fonttitle=\normalfont\large,
  overlay={
    \draw[black, ultra thick] 
%    \draw[gblue700, ultra thick] 
    ([xshift=6pt,yshift=-0.5\pgflinewidth]title.south west) --
    ([xshift=-6pt,yshift=-0.5\pgflinewidth]title.south east);},
  left=1mm,
  right=1mm,
  top=1mm,
  bottom=1mm,
  middle=1mm,
  fuzzy shadow={0.6ex}{-0.6ex}{1pt}{0.5pt}{black!80!white},
  #1,title=#2}
\makeatletter
\tikzset{beamer@color/.style={top color=white,bottom color=white}}
\makeatother

}{

\iftoggle{chapter}{%

	\declaretheorem[style=plain,qed=$\lrcorner$,numberwithin=chapter]{theorem}

}{

	\declaretheorem[style=plain,qed=$\lrcorner$,numbered=yes]{theorem}

}


% Theorem environments

\declaretheorem[style=definition,qed=$\lrcorner$,sibling=theorem]{lemma}
\declaretheorem[style=definition,qed=$\lrcorner$,sibling=theorem]{corollary}
\declaretheorem[style=definition,qed=$\lrcorner$,sibling=theorem]{proposition}
\declaretheorem[style=definition,qed=$\lrcorner$,sibling=theorem]{claim}
\declaretheorem[style=definition,qed=$\lrcorner$,sibling=theorem]{conjecture}

% Remark and definition environments
\declaretheorem[style=definition,qed=$\lrcorner$,sibling=theorem]{assumption}
\crefname{assumption}{Assumption}{Assumptions}
\Crefname{assumption}{Assumption}{Assumptions}
\declaretheorem[style=definition,qed=$\lrcorner$,sibling=theorem]{definition}
\declaretheorem[style=definition,qed=$\lrcorner$,sibling=theorem]{remark}
\declaretheorem[style=definition,qed=$\lrcorner$,sibling=theorem]{problem}
\crefname{problem}{Problem}{Problems}
\Crefname{problem}{Problem}{Problems}

\declaretheorem[style=definition,qed=$\lrcorner$,numberwithin=section]{question}
\declaretheorem[style=definition,qed=$\lrcorner$,numbered=no]{answer}

% Example environments

\declaretheorem[style=definition,qed=$\lrcorner$,sibling=theorem]{example}

% Proof environment
\declaretheoremstyle[
  spaceabove=6pt, spacebelow=6pt,
  headfont=\normalfont\slshape,
  notefont=\normalfont\slshape, notebraces={(}{)},
  bodyfont=\normalfont,
  postheadspace=1em,
  qed=$\blacksquare$
]{proofstyle}
\declaretheorem[name=Proof,style=proofstyle,numbered=no]{theproof}

}

%\usepackage{fixmath}

%%%%%%%%%%%%%%%%%%%%
% BIBLATEX PACKAGE %
%%%%%%%%%%%%%%%%%%%%

% To be loaded locally after the language package

\usepackage[natbib=true,backend=biber,style=alphabetic,sorting=anyvt,
firstinits=true,doi=false,isbn=false,url=false,backref=true,texencoding=utf8,bibencoding=utf8]{biblatex}

\IfFileExists{./../Library.bib}{

 \addbibresource{./../../../Library}
 
}{

  %Must be used with biber as backend
  %Execute with "biber SeparableMetrics"
  \addbibresource[location=remote]{https://www.dropbox.com/s/r00o9lw76cksxvj/Library.bib?dl=1} 
}






\graphicspath{ {./imgs/} }
\newcommand{\svginput}[1]{\input{imgs/#1}}

%%%%%%%%%%%%%
% TITLEPAGE %
%%%%%%%%%%%%%

% Needs to be loaded locally

\usetikzlibrary{shapes,arrows}
\setbeamerfont{author}{size=\LARGE}
\setbeamerfont{institute}{size=\normalsize\itshape}
\setbeamerfont{title}{size=\fontsize{24}{30}\bfseries}
\setbeamerfont{subtitle}{size=\Large\normalfont\slshape}

\setbeamertemplate{title page}{%
\begin{tikzpicture}[remember picture,overlay]
\fill[gblue700]
  ([yshift=0pt]current page.west) rectangle ([yshift=-\headheight] current page.north east);
\node[anchor=east] 
  at ([yshift=85pt]current page.south east) (author)
  {\parbox[t]{\paperwidth}{\centering%
    \usebeamerfont{author}\textcolor{gblue700}{%
    \textpdfrender{
    TextRenderingMode=FillStroke,
    FillColor=gblue700,
    LineWidth=.1ex,
    }{\insertauthor}}}};
%\node[anchor=south east] 
%  at ([yshift=0pt]current page.south east) (institute)
%  {\parbox[t]{.78\paperwidth}{\raggedleft%
%    \usebeamerfont{institute}\textcolor{black}{\insertinstitute}}};
\node[anchor=south] 
  at ([yshift=0pt]current page.south) (logo)
  {\parbox[t]{\paperwidth}{\centering%
    \usebeamercolor[fg]{titlegraphic}\inserttitlegraphic}};
\node[anchor=center]
  at ([yshift=50pt,xshift=0pt]current page.center) (title)
  {\parbox[t]{\textwidth}{\centering%
 \usebeamerfont{author}\textcolor{white}{%
    \textpdfrender{
    TextRenderingMode=FillStroke,
    FillColor=white,
    LineWidth=.1ex,
    }{\inserttitle}}}};
\node[anchor=east]
  at ([yshift=-60pt,xshift=-20pt]current page.east) (subtitle)
  {\parbox[t]{.6\paperwidth}{\raggedleft\usebeamerfont{subtitle}\textcolor{black}{\insertsubtitle}}};
\end{tikzpicture}
}


\title[XeLaTeX Template]{My XeLaTeX Template}
\author[humberto.shiromoto@ieee.org]{Humberto \sc Stein Shiromoto}
\institute{The University of Sydney}
\date[19/05/15]{\today}
\titlegraphic{
\begin{minipage}[b]{0.95\linewidth}
\begin{figure}[htbp!]
\centering
\includegraphics[width=2.5cm]{./imgs/Usyd_new_logo}
\end{figure}
\end{minipage}
}

\begin{document}

\tikzstyle{every picture}+=[remember picture]

% By default all math in TikZ nodes are set in inline mode. Change this to
% displaystyle so that we don't get small fractions.
\everymath{\displaystyle}
	
	{
	%Highlight section in navigation bar
	\setbeamertemplate{section in head/foot shaded}[default][100]
	\setbeamercolor{upper separation line head}{bg=gblue700}
	% Highlight dots
	\setbeamertemplate{mini frame in other subsection}[default][100]
	% Highlight and fill dots
%	\setbeamertemplate{mini frame in other subsection}{%
%    \begin{pgfpicture}{0pt}{0pt}{0.1cm}{0.1cm}%
%        \pgfpathcircle{\pgfpoint{0.05cm}{0.05cm}}{0.05cm}%
%        \pgfusepath{fill,stroke}%
%    \end{pgfpicture}%
%    }%
		
	
	
	\frame{
		\maketitle	
	}	}
	
	\section[Text, items, equations, and figures]{Tests with text, items, equations, and figures}	
	\subsection[First subsection]{Text and Equations}	
	\frame[c]{
		\frametitle{The first frame: text}
		\begin{adjustwidth}{-2em}{-2em}
		\blindtext
		\end{adjustwidth}
	}

	\frame[c]{
		\frametitle{The First Frame: equation}
		\begin{adjustwidth}{-2em}{-2em}
		% \blindmathpaper % Gives errors!
		\end{adjustwidth}
	}

	\subsection{Items and lists}
	\frame[c]{
		\frametitle{The second frame: items}
		\begin{adjustwidth}{-2em}{-2em}
		\blindlist{itemize}[3]	
		\end{adjustwidth}
	}

	\frame[c]{
		\frametitle{The third frame: numbered items}
		\begin{adjustwidth}{-2em}{-2em}
		\blindlist{enumerate}[5]	
		\end{adjustwidth}
	}
	
	\frame[c]{
		\frametitle{The fourth frame: description list}
		\begin{adjustwidth}{-2em}{-2em}
		\blindlist{description}[2]	
		\end{adjustwidth}
	}
	
	\subsection{Figure}
	\frame[c]{
		\frametitle{The Second Frame: a figure}
		\begin{adjustwidth}{-2em}{-2em}
		\begin{figure}[htpb!]
			\centering
			\scalebox{0.75}{\input{./imgs/set_evolution.eps_tex}}
		\end{figure}
		\end{adjustwidth}
		}
	\section{Tests with boxes}
	\subsection{Box without shadows}
	\frame[c]{
		\frametitle{Without shadows}
		
		\begin{tcolorbox}[title=A test]
This is an equation
\begin{equation*}
	\int_{X(t,\mathbf{Z})}\rho\,dx-\int_\mathbf{Z}=\int_0^t\int_{X(s,\mathbf{Z})}\nabla(\rho f)(x)\,dx
\end{equation*}
\end{tcolorbox}\par\bigskip
\begin{tcolorbox}
This is a tcolorbox.
\end{tcolorbox}\par\bigskip
}

	\subsection{Lifted shadows}
	
	\frame[c]{
	\frametitle{Lifted shadows}
				\begin{tcolorbox}[drop lifted shadow]
This is a tcolorbox.
\end{tcolorbox}\par\bigskip
\begin{tcolorbox}[title=Another shadow,
drop lifted shadow]
This is a tcolorbox.
\end{tcolorbox}\par\bigskip
\begin{tcolorbox}[drop lifted shadow]
\begin{equation*}
% Todo: check why \dot{x} doesn't work
	\dfrac{dx}{dt}=f(x)
\end{equation*}
\end{tcolorbox}\par\bigskip
\begin{tcolorbox}[code={\pgfkeysalsofrom\redblock},title=A test,drop lifted shadow]
Test.
\end{tcolorbox}		
	}
	
	\subsection{Midday Shadows}
	
	\frame[c]{
		\frametitle{Midday shadows}
		
		\begin{tcolorbox}[drop fuzzy midday shadow,
enhanced,colback=red!5!white,colframe=gred900]
This is a tcolorbox.
\end{tcolorbox}\par\bigskip
\begin{tcolorbox}[title=Another shadow,drop fuzzy midday shadow,
enhanced,colback=red!5!white,colframe=gred900]
This is a tcolorbox.
\end{tcolorbox}
}

% \frame[t]{
% 	% Requires definition of \normal and \light fonts
% 	\frametitle{Variations on the title: available if \texttt{normal} and \texttt{light} defined}
% 	\begin{tcolorbox}[title=Another shadow,
% drop lifted shadow]
% This is a tcolorbox.
% \end{tcolorbox}\par\bigskip
% \begin{tcolorbox}[fonttitle=\normal,title=Another shadow,
% drop lifted shadow]
% This is a tcolorbox.
% \end{tcolorbox}\par\bigskip
% \begin{tcolorbox}[fonttitle=\light,title=Another shadow,
% drop lifted shadow]
% This is a tcolorbox.
% \end{tcolorbox}
% }

			\frame[c]{
	\frametitle{Small boxes with midday, and lifted shadows}
	\begin{minipage}[c]{0.30\linewidth}
	\begin{tcolorbox}[code={\pgfkeysalsofrom\postitblock},drop fuzzy midday shadow,hbox]
Test.
\end{tcolorbox}
\begin{tcolorbox}[code={\pgfkeysalsofrom\whiteblock},drop fuzzy midday shadow,hbox]
Test.
\end{tcolorbox}
	\end{minipage}
	\begin{minipage}[c]{0.30\linewidth}
	\begin{tcolorbox}[code={\pgfkeysalsofrom\postitblock},drop lifted shadow,hbox]
Test.
\end{tcolorbox}
\begin{tcolorbox}[code={\pgfkeysalsofrom\whiteblock},drop lifted shadow,hbox]
Test.
\end{tcolorbox}
	\end{minipage}
	\begin{minipage}[c]{0.30\linewidth}
	\begin{tcolorbox}[code={\pgfkeysalsofrom\postitblock},drop fuzzy shadow south,hbox]
Test.
\end{tcolorbox}
\begin{tcolorbox}[code={\pgfkeysalsofrom\whiteblock},drop fuzzy shadow south,hbox]
Test.
\end{tcolorbox}
	\end{minipage}
	
	\begin{tcolorbox}[code={\pgfkeysalsofrom\postitblock},drop lifted shadow]
	\begin{varwidth}{\textwidth}
	The text
	
       $\int_{X(t,\mathbf{Z})}\rho\,dx-\int_\mathbf{Z}=\int_0^t\int_{X(s,\mathbf{Z})}\nabla(\rho f)(x)\,dx$
       \end{varwidth}
       \end{tcolorbox}
	}
	
	\section{Theorem-like environments}
	\subsection{Theorem environments}
	\frame[c]{	
		\frametitle{ToDo}
	
	}
	\subsection{Definition environments}
	
	\subsection{Example environments}
		
\end{document}
